%! Exponential Growth and Decay — Unit 7, Lesson 7.1 — Group Activity (Answer Key)
\documentclass[10pt]{article}
\usepackage{algebra2-article}
\usepackage{algebra2-key}   % pulls in -boxes; do NOT also load -boxes
\newcommand{\TallMath}[1]{$\displaystyle #1\rule[-1.4em]{0pt}{3.2em}$}
\newcommand{\lgrid}[4]{%
  \draw[step=1, linegray!40, very thin] (#1,#3) grid (#2,#4);
  \draw[->, semithick, charcoal] (#1,0)--(#2,0) node[below right, font=\scriptsize]{$x$};
  \draw[->, semithick, charcoal] (0,#3)--(0,#4) node[left, font=\scriptsize]{$y$};}
% #2 is ln(b):  ln2 = 0.693147   ln3 = 1.098612
\newcommand{\expcurve}[4]{\draw[very thick, forest, domain=#3:#4, samples=120, smooth]
  plot (\x,{#1*exp((#2)*(\x))});}

\begin{document}

\pageheader{Unit 7, Lesson 7.1}{Group Activity --- Answer Key}
\namepartnerperiod

\vspace{0.08in}

\begin{headlinebox}{forestbg}
\small\textbf{Building the Model.} Yesterday the equation was handed to you. Today you write it. Every
problem below gives you a \emph{source} --- a table, a graph, or a sentence --- and asks for the same
two numbers: the initial value $a$ and the factor $b$. With your partner, build each model, then
\textbf{check it}: does $x=0$ give back the starting amount, and is $b$ on the correct side of $1$?
\end{headlinebox}

\vspace{0.06in}

\begin{tcolorbox}[colback=white, colframe=black!40, title=\textbf{Tier R --- Remediate},
    fonttitle=\bfseries, arc=1.5mm, left=3mm, right=3mm, top=2mm, bottom=2mm, breakable]
\textbf{Part 1 --- Percent to factor.} Fill in each row. Remember: $b=1+r$ when it grows,
$b=1-r$ when it shrinks.
\begin{center}
\renewcommand{\arraystretch}{1.55}
\begin{tabularx}{\linewidth}{@{}X l l X@{}}
\hline
\textbf{The sentence says\dots} & \textbf{Rate $r$} & \textbf{Grows or shrinks?} & \textbf{Factor $b$} \\
\hline
increases by $6\%$ each year & \ans{0.06} & \ans{grows} & \ans{$1.06$} \\
decreases by $6\%$ each year & \ans{0.06} & \ans{shrinks} & \ans{$0.94$} \\
grows by $25\%$ each week & \ans{0.25} & \ans{grows} & \ans{$1.25$} \\
loses $40\%$ of it each hour & \ans{0.40} & \ans{shrinks} & \ans{$0.60$} \\
\hline
\end{tabularx}
\end{center}
The first two rows describe the same percent. Why do they give \emph{different} bases?
\ans{One \emph{adds} the $6\%$ to what you already have ($1+0.06$); the other \emph{takes it away} ($1-0.06$).}

\vspace{6pt}
\textbf{Part 2 --- One full model.} An ant colony has $60$ ants and grows by $25\%$ each week.
\begin{enumerate}[label=(\alph*), leftmargin=1.9em, itemsep=6pt]
  \item $a=$ \ans{60} \quad $r=$ \ans{0.25} \quad $b=$ \ans{1.25}
  \item Write the model: $N(w)=$ \ans{$60(1.25)^{w}$} \quad (where $w$ is weeks)
  \item \textbf{Check it.} What does your model give at $w=0$? \ans{60} \;
        Does that match the story? \ans{yes}
  \item After $1$ week: $N(1)=$ \ans{75}. \; Is that $25\%$ more than $60$?
        \ans{yes ($60+15$)}
  \item After $4$ weeks: $N(4)=$ \ans{$60(1.25)^{4}=146.48\ldots$} $\approx$ \ans{146} ants. \\
        (Ants come in whole numbers --- round down.)
  \item A classmate wrote $N(w)=60(0.25)^{w}$. What does \emph{that} model describe?
        \ans{A colony where only $25\%$ survives each week --- a $75\%$ \emph{loss}, not a $25\%$ gain.}
\end{enumerate}
\end{tcolorbox}

\begin{tcolorbox}[colback=white, colframe=black!40, title=\textbf{Tier A --- Approaching Proficiency},
    fonttitle=\bfseries, arc=1.5mm, left=3mm, right=3mm, top=2mm, bottom=2mm, breakable]
\textbf{Three sources, one form.} Build $y=ab^{x}$ from each.

\vspace{5pt}
\textbf{(a) From a table.}
\begin{center}
\renewcommand{\arraystretch}{1.25}
\begin{tabular}{|c|c|c|c|c|}
\hline
$x$ & $0$ & $1$ & $2$ & $3$ \\ \hline
$y$ & $9$ & $27$ & $81$ & $243$ \\ \hline
\end{tabular}
\end{center}
Ratios: $\frac{27}{9}=$ \ans{3}, \; $\frac{81}{27}=$ \ans{3}. \quad
$a=$ \ans{9} \quad $b=$ \ans{3} \quad $y=$ \ans{$9\cdot 3^{x}$}

\vspace{6pt}
\textbf{(b) From a graph.} Use the three labeled points.
\begin{center}
\begin{minipage}[c]{0.44\linewidth}
\centering
\begin{tikzpicture}[scale=0.44]
  \lgrid{-1}{5}{-1}{10}
  \begin{scope}
    \clip (-1,-1) rectangle (5,10);
    \expcurve{9}{-1.098612}{-0.09}{5}
  \end{scope}
  \draw[navy, dashed, semithick] (-1,0)--(5,0);
  \fill[charcoal] (0,9) circle (3pt) node[right, font=\scriptsize]{$(0,9)$};
  \fill[charcoal] (1,3) circle (3pt) node[right, font=\scriptsize]{$(1,3)$};
  \fill[charcoal] (2,1) circle (3pt) node[above right, font=\scriptsize]{$(2,1)$};
\end{tikzpicture}
\end{minipage}
\begin{minipage}[c]{0.52\linewidth}
$a=$ \ans{9} \; (read it off the $y$-axis) \\[6pt]
$b=\dfrac{3}{9}=$ \ans{$\frac13$} \\[6pt]
Growth or decay? \ans{decay} \\[6pt]
$y=$ \ans{$9\left(\frac13\right)^{x}$} \\[6pt]
Horizontal asymptote: \ans{$y=0$}
\end{minipage}
\end{center}

\vspace{4pt}
\textbf{(c) From a story.} A savings account holds \$$800$ and earns $4\%$ interest each year.
\\[3pt]
$a=$ \ans{800} \quad $r=$ \ans{0.04} \quad $b=$ \ans{1.04} \quad
$S(t)=$ \ans{$800(1.04)^{t}$} \\[3pt]
$S(6)=$ \ans{$800(1.2653)$} $\approx$ \ans{\$$1012.26$}. \; What does that number mean?
\ans{After $6$ years the account holds about \$$1012.26$ --- about \$$212$ of interest.}

\vspace{6pt}
\textbf{Justify.} A deer population starts at $500$ and \textbf{declines $20\%$ each year}. Three
students hand in three different models:
\begin{center}
\small
\renewcommand{\arraystretch}{1.3}
\begin{tabularx}{0.94\linewidth}{@{}l l X@{}}
\hline
Maya & $P(t)=500(1.2)^{t}$ & \ans{models a $20\%$ \emph{increase}} \\
Devon & $P(t)=500(0.2)^{t}$ & \ans{models an $80\%$ \emph{decrease}} \\
Rosa & $P(t)=500(0.8)^{t}$ & \ans{correct} \\
\hline
\end{tabularx}
\end{center}
Who is right? Then --- and this is the real question --- say what each \emph{wrong} model actually
describes. Be specific: give the percent and the direction.
\ansline{\textbf{Rosa} is right: a $20\%$ decline leaves $80\%$, so $b=1-0.20=0.8$.\\[3pt]
\textbf{Maya} used $b=1.2$. That base is greater than $1$, so her herd \emph{grows} by $20\%$ a
year --- she added the rate instead of subtracting it.\\[3pt]
\textbf{Devon} used $b=0.2$, the percent itself. That says only $20\%$ of the deer are left each
year, which is an $80\%$ decline --- far steeper than the story. After $3$ years Devon predicts
$4$ deer; Rosa predicts $256$.}
\end{tcolorbox}

\begin{tcolorbox}[colback=white, colframe=black!40, title=\textbf{Tier E --- Extension},
    fonttitle=\bfseries, arc=1.5mm, left=3mm, right=3mm, top=2mm, bottom=2mm, breakable]
\textbf{Part 1 --- Two points, neither of them the $y$-intercept.} Divide to cancel $a$, then solve
for $b$, then go back for $a$.
\begin{enumerate}[label=(\alph*), leftmargin=1.9em, itemsep=7pt]
  \item The curve $y=ab^{x}$ passes through $(2,45)$ and $(5,1215)$. \\[3pt]
        $\dfrac{1215}{45}=$ \ans{27} $=b^{\,\square}$ where $\square=$ \ans{3}, so
        $b=$ \ans{3}. \\[3pt]
        Now use $(2,45)$: \; $45=a\cdot$ \ans{9}, so $a=$ \ans{5}. \quad
        $y=$ \ans{$5\cdot 3^{x}$}
  \item The curve $y=ab^{x}$ passes through $(1,60)$ and $(4,7.5)$. \\[3pt]
        $b=$ \ans{$\frac12$} \quad $a=$ \ans{120} \quad $y=$ \ans{$120\left(\frac12\right)^{x}$} \\[3pt]
        How could you tell this one would be \emph{decay} before doing any arithmetic?
        \ans{The output got \emph{smaller} as $x$ got larger, so $b$ has to be between $0$ and $1$.}
\end{enumerate}

\vspace{5pt}
\textbf{Part 2 --- Run the translation backwards.} Each model is already built. Report the starting
amount and the percent change \emph{per period}, and say which direction.
\begin{center}
\renewcommand{\arraystretch}{1.55}
\begin{tabularx}{\linewidth}{@{}l X X X@{}}
\hline
\textbf{Model} & \textbf{Initial value} & \textbf{Percent rate $r$} & \textbf{Increase or decrease?} \\
\hline
$y=500(1.045)^{t}$ & \ans{500} & \ans{$4.5\%$} & \ans{increase} \\
$y=2200(0.91)^{t}$ & \ans{2200} & \ans{$9\%$} & \ans{decrease} \\
$y=75(2)^{t}$ & \ans{75} & \ans{$100\%$} & \ans{increase (it doubles)} \\
$y=1000(1.005)^{t}$ & \ans{1000} & \ans{$0.5\%$} & \ans{increase} \\
\hline
\end{tabularx}
\end{center}

\vspace{4pt}
\textbf{Part 3 --- The same first step, two different futures.} A tank holds $240$ gallons. Compare
two leaks:
\begin{center}
\small
\renewcommand{\arraystretch}{1.3}
\begin{tabularx}{0.94\linewidth}{@{}l X@{}}
\hline
\textbf{Leak A} & The tank loses $60$ gallons every day. \\
\textbf{Leak B} & The tank loses $25\%$ of whatever is in it every day. \\
\hline
\end{tabularx}
\end{center}
\begin{enumerate}[label=(\alph*), leftmargin=1.9em, itemsep=6pt]
  \item Write both models. \quad $A(d)=$ \ans{$240-60d$} \quad $B(d)=$ \ans{$240(0.75)^{d}$}
  \item Fill in the table (round to the nearest hundredth).
    \begin{center}
    \renewcommand{\arraystretch}{1.3}
    \begin{tabular}{|c|c|c|c|c|c|c|}
    \hline
    Day $d$ & $0$ & $1$ & $2$ & $3$ & $4$ & $5$ \\ \hline
    $A(d)$ & $240$ & \ans{180} & \ans{120} & \ans{60} & \ans{0} & \ans{$-60$} \\ \hline
    $B(d)$ & $240$ & \ans{180} & \ans{135} & \ans{101.25} & \ans{75.94} & \ans{56.95} \\ \hline
    \end{tabular}
    \end{center}
  \item On day $1$ the two leaks lose the \emph{same} amount. Show that, then explain why they come
        apart after that.
        \ansline{$25\%$ of $240$ is $60$, so both tanks drop to $180$ on day $1$. After that Leak A
        keeps subtracting the \emph{same} $60$ gallons, while Leak B takes $25\%$ of a
        \emph{smaller} amount each time --- $45$, then $33.75$, then $25.31$. Constant difference
        versus constant ratio.}
  \item Leak A empties the tank. On what day? \ans{day 4} \; Leak B never does.
        Use Lesson 7.0's language --- \textbf{asymptote} --- to explain why not.
        \ansline{$B(d)=240(0.75)^{d}$ has the horizontal asymptote $B=0$. Multiplying a positive
        number by $0.75$ always leaves a positive number, so the amount gets arbitrarily close to
        $0$ but never reaches it --- $0$ is a boundary of the range, not a value the model attains.}
  \item After two weeks ($d=14$), Leak A's formula gives \ans{$-600$} gallons and Leak B's
        gives about \ans{$4.28$} gallons. One of those answers is \emph{impossible} and the
        other is merely \emph{unrealistic}. Explain both.
        \ansline{Leak A's $-600$ is \emph{impossible}: a tank cannot hold a negative amount of
        water. The linear model is only valid until day $4$, when the tank runs dry --- past that
        its contextual domain has ended.\\[3pt]
        Leak B's $4.28$ gallons is a real, possible number, but the model claims the tank still has
        \emph{something} in it forever, no matter how many days pass. That is the asymptote talking.
        In reality the last few drops evaporate or drain, so the model is useful early and
        meaningless in the far tail.}
\end{enumerate}
\end{tcolorbox}

\begin{teachernote}
\textbf{Tier R Part 1's follow-up question} (``same percent, different bases'') is the whole point of
the tier --- do not let students skip it to get to Part 2. If a group writes $b=0.06$ anywhere, send
them back to Warm-Up item 1 rather than correcting it for them.

\textbf{Tier A's Justify is the assessment item of the activity.} Identifying Rosa is easy; the
work is diagnosing Maya and Devon. Maya's error is a \emph{sign} error ($1+r$ instead of $1-r$) and
Devon's is the \emph{percent-as-base} error --- two different misconceptions that need two different
re-teaches, which is exactly why the question asks students to describe each wrong model rather than
just cross it out. The numbers at $t=3$ ($4$ deer versus $256$) are worth putting on the board.

\textbf{Tier E Part 1} is the only place two-point building gets real practice, and (b) is harder than
(a) because $b$ comes out a fraction. Groups that stall usually forget that $\frac{7.5}{60}=b^{3}$
requires a \emph{cube} root, not a square root --- the exponent is the \emph{gap} between the
$x$-values, $4-1=3$.

\textbf{Tier E Part 3} is the linear-versus-exponential contrast from Lesson 7.0's fingerprint test,
now in context, and it sets up two later lessons at once: the contextual-domain question returns in
7.4's modeling, and part (e) is a rehearsal for 7.5's ``extrapolation breaks'' beat. Both tanks
losing exactly $60$ gallons on day $1$ is deliberate --- it kills the idea that ``exponential just
means faster.''

\textbf{Pacing:} Tier R is one page of work and should take about ten minutes; Tier A about fifteen.
Tier E is genuinely long --- assign Part 1 and \emph{either} Part 2 or Part 3 if the period is tight.
\end{teachernote}

\end{document}
